% META: {"id": "TSPE-GEOMETRIE-ESPACE-QCM", "chapitre": "TSPE-GEOMETRIE-ESPACE", "type_objet": "qcm", "genere_depuis": "chapitres/TSPE-GEOMETRIE-ESPACE/qcm/TSPE-GEOMETRIE-ESPACE-QCM.json", "status": "generated"}
% Fichier genere par scripts/build_qcm_tex.py — ne pas editer a la main.

\section*{\textcolor{chapcolor}{\MakeUppercase{QCM d'auto-evaluation — Géométrie dans l'espace}}}

\begin{center}
\textit{Pour chaque question, une seule reponse est exacte.}
\end{center}

\begin{enumerate}

\item[] \textbf{Capacite C3}

\item \textbf{[Q1]} Deux droites de l'espace non paralleles sont :
  \begin{enumerate}[label=\Alph*.]
    \item toujours secantes.
    \item toujours non coplanaires.
    \item secantes ou non coplanaires.
    \item toujours confondues.
  \end{enumerate}

\item[] \textbf{Capacite C7}

\item \textbf{[Q2]} Deux vecteurs $\vec{u}$ et $\vec{v}$ sont orthogonaux si et seulement si :
  \begin{enumerate}[label=\Alph*.]
    \item $\vec{u}\cdot\vec{v}=1$.
    \item $\vec{u}\cdot\vec{v}=0$.
    \item $\vec{u}=\vec{v}$.
    \item $\|\vec{u}\|=\|\vec{v}\|$.
  \end{enumerate}

\item[] \textbf{Capacite C13}

\item \textbf{[Q3]} Le vecteur normal au plan d'equation $2x-3y+z+1=0$ est :
  \begin{enumerate}[label=\Alph*.]
    \item $(2,-3,1)$.
    \item $(2,3,1)$.
    \item $(1,-3,2)$.
    \item $(-2,3,1)$.
  \end{enumerate}

\item[] \textbf{Capacite C12}

\item \textbf{[Q4]} Une representation parametrique de droite comporte :
  \begin{enumerate}[label=\Alph*.]
    \item un seul parametre.
    \item deux parametres.
    \item trois parametres.
    \item aucun parametre.
  \end{enumerate}

\item[] \textbf{Capacite C11}

\item \textbf{[Q5]} Le projete orthogonal $H$ d'un point $M$ sur un plan $\mathcal{P}$ verifie :
  \begin{enumerate}[label=\Alph*.]
    \item $H$ est le point de $\mathcal{P}$ le plus eloigne de $M$.
    \item $H$ est le point de $\mathcal{P}$ le plus proche de $M$.
    \item $H$ est un point quelconque de $\mathcal{P}$.
    \item $H=M$ toujours.
  \end{enumerate}

\end{enumerate}

% NEXUS-QCM-TEACHER-ONLY-BEGIN
\ifnxVersionProfesseur
\clearpage
\section*{Cle de correction — reservee au professeur}

\begin{center}
\begin{tabular}{lll}
\hline
Question & Capacite & Reponse exacte \\
\hline
Q1 & C3 & \textbf{C} \\
Q2 & C7 & \textbf{B} \\
Q3 & C13 & \textbf{A} \\
Q4 & C12 & \textbf{A} \\
Q5 & C11 & \textbf{B} \\
\hline
\end{tabular}
\end{center}

\subsection*{Diagnostic des reponses erronees}

\begin{itemize}
  \item \textbf{Q1}
  \begin{itemize}
    \item \textbf{A} — Dans l'espace, deux droites non paralleles peuvent etre gauches : elles ne sont alors pas coplanaires et ne se coupent pas. \emph{Renvoi : M1.}
    \item \textbf{B} — Deux droites non paralleles peuvent etre coplanaires et secantes ; elles ne sont donc pas toujours non coplanaires. \emph{Renvoi : M1.}
    \item \textbf{D} — Des droites confondues ont la meme direction et sont donc paralleles ; cette possibilite est exclue par l'hypothese. \emph{Renvoi : M1.}
  \end{itemize}
  \item \textbf{Q2}
  \begin{itemize}
    \item \textbf{A} — Un produit scalaire egal a $1$ n'impose pas un angle droit ; l'orthogonalite est caracterisee par un produit scalaire nul. \emph{Renvoi : M1.}
    \item \textbf{C} — Si $\vec u=\vec v$ et ce vecteur est non nul, alors $\vec u\cdot\vec v=\|\vec u\|^2>0$ : ils ne sont pas orthogonaux. \emph{Renvoi : M1.}
    \item \textbf{D} — L'egalite des normes fixe seulement les longueurs ; deux vecteurs de meme norme peuvent former n'importe quel angle. \emph{Renvoi : M1.}
  \end{itemize}
  \item \textbf{Q3}
  \begin{itemize}
    \item \textbf{B} — Le coefficient de $y$ dans l'equation est $-3$ ; remplacer ce coefficient par $3$ ne donne pas un vecteur normal au plan. \emph{Renvoi : M1.}
    \item \textbf{C} — Les composantes d'un vecteur normal sont les coefficients de $x$, $y$ et $z$ dans cet ordre ; ici $2$, $-3$ et $1$. \emph{Renvoi : M1.}
    \item \textbf{D} — Ce vecteur change les signes des deux premieres composantes seulement ; il n'est donc pas proportionnel a $(2,-3,1)$. \emph{Renvoi : M1.}
  \end{itemize}
  \item \textbf{Q4}
  \begin{itemize}
    \item \textbf{B} — Deux parametres independants engendrent en general une surface plane, alors qu'une droite ne possede qu'une direction libre. \emph{Renvoi : M1.}
    \item \textbf{C} — Trois parametres independants peuvent parcourir l'espace entier ; une droite s'ecrit avec un seul parametre reel. \emph{Renvoi : M1.}
    \item \textbf{D} — Sans parametre, on ne decrit pas le deplacement depuis un point selon un vecteur directeur ; la representation ne serait pas parametrique. \emph{Renvoi : M1.}
  \end{itemize}
  \item \textbf{Q5}
  \begin{itemize}
    \item \textbf{A} — Un plan est non borne, donc il ne possede pas de point le plus eloigne de $M$ ; la projection realise au contraire la distance minimale. \emph{Renvoi : M1.}
    \item \textbf{C} — Le projete n'est pas quelconque : c'est le pied de la perpendiculaire issue de $M$, ce qui assure la distance minimale au plan. \emph{Renvoi : M1.}
    \item \textbf{D} — On a $H=M$ seulement lorsque $M$ appartient deja au plan ; sinon le projete $H$ est un point distinct situe sur le plan. \emph{Renvoi : M1.}
  \end{itemize}
\end{itemize}
\fi
% NEXUS-QCM-TEACHER-ONLY-END
